\section{The Analytical Shift: From Monitoring to Observability}
\label{sec:analytical-shift}

According to recent literature \cite{kosinska_2023}, observability is often defined as consisting of core pillars: metrics, logging, and tracing, with modern frameworks also encompassing events \cite{bhatia_2025}. However, as cautioned by Hartikainen (2024), the mere fact that a system possesses these pillars does not automatically make it observable. The value lies in the semantic conversion of these disjointed signals into actionable insights.

This distinction necessitates a critical boundary between traditional monitoring and modern observability. Although often used interchangeably in the industry, they represent different operational paradigms. Monitoring tracks a system's health using a predefined set of metrics to detect known failures (the "known unknowns"). It is essentially the "alarm" warning of an existing problem. 

Observability, in contrast, is the analytical capability to investigate and determine \textit{where} and \textit{why} a failure occurred, especially for unforeseen issues (the "unknown unknowns") in dynamic, distributed systems. While monitoring tells the operator that a system is broken, observability provides the mathematical and architectural means to infer the broken internal state, fulfilling Kalman's original premise of system controllability. Therefore, monitoring is a prerequisite for observability, but relying solely on monitoring in cloud-native architectures invariably leads to an inability to diagnose emergent failures.

Evaluating the current landscape through Kalman's theoretical lens reveals a critical shortcoming in modern cloud-native environments. Existing solutions have vastly improved the extraction of ``external outputs'' (the telemetry data via robust tools like Prometheus or Jaeger), but they often fall short of achieving true observability. The literature demonstrates that merely possessing isolated dashboards for metrics, logs, and traces creates an advanced monitoring setup, not an observable system \cite{usman_2022}. Because these tools lack native semantic correlation, they fail to seamlessly reconstruct the ``internal state'' of the application without manual human intervention.

This cognitive burden during Root Cause Analysis (RCA), exacerbated by cloud-native heterogeneity and the sheer volume of disjointed multimodal data, is precisely where the current state-of-the-art falls short of Kalman's ideal \cite{han_2024, gomes_2025, li_2024}. As highlighted by Han et al. (2024), the lack of a holistic approach to correlate this heterogeneous data inherently limits generalizability and increases the complexity of decision-making for engineers. This reinforces the necessity for a unified reference architecture that mathematically and operationally links these disjointed signals.

In the next chapter, the research roadmap of observability of cloud native applications will be revealed. Rapid Review approach was chosen as it aims to swiftly synthesize scientific evidence on the rapidly evolving field of observability in cloud-native applications without compromising systematic rigor.